\documentclass[12pt,letterpaper]{article}

%Packages
\usepackage{amsmath}
\usepackage{amsthm}
\usepackage{amsfonts}
\usepackage{amssymb}
\usepackage{amscd}
\usepackage{dsfont}
\usepackage{enumerate}
\usepackage{fancyhdr}
\usepackage{mathrsfs}
\usepackage{bbm}
\usepackage{framed}
\usepackage{mdframed}
\usepackage{cancel}
\usepackage{float}
\usepackage{mathtools}
%\usepackage[]{mcode}
\usepackage{graphicx}
\usepackage{tikz}
\usepackage{hyperref}

%Page formatting
\usepackage[letterpaper,voffset=-.5in,bmargin=3cm,footskip=1cm]{geometry}
\setlength{\parindent}{0.0in}
\setlength{\parskip}{0.1in}
\allowdisplaybreaks
\headheight 15pt
\headsep 10pt

% Common Commands
\newcommand\N{\mathbb N}
\newcommand\Z{\mathbb Z}
\newcommand\R{\mathbb R}
\newcommand\Q{\mathbb Q}
\newcommand\lcm{\operatorname{lcm}}
\newcommand\setbuilder[2]{\ensuremath{\left\{#1\;\middle|\;#2\right\}}}
\newcommand\E{\operatorname{E}}
\newcommand\V{\operatorname{V}}
\newcommand\Pow{\ensuremath{\operatorname{\mathcal{P}}}}

\DeclarePairedDelimiter\ceil{\lceil}{\rceil}
\DeclarePairedDelimiter\floor{\lfloor}{\rfloor}

% 22 style
\newcommand\hint[1]{\textbf{Hint}: #1}
\newcommand\note[1]{\textbf{Note}: #1}
\newenvironment{22enumerate}{\begin{enumerate}[a.]\itemsep0em}{\end{enumerate}}
\newenvironment{22itemize}{\begin{itemize}\itemsep0em}{\end{itemize}}
\fancypagestyle{firstpagestyle} {
  \renewcommand{\headrulewidth}{0pt}%
  \lhead{\textbf{CSCI 0220}}%
  \chead{\textbf{Discrete Structures and Probability}}%
  \rhead{Littman}%
}

\pagestyle{fancyplain}

\newcommand\EX{\mathds{E}}
\newcommand\PR{\mathds{P}}
\newcommand\zlam{Z_\lambda}
\DeclareMathOperator*{\argmin}{arg\,min}

%%%%%%%%%%%%%%%% COMPILE WITH \soltrue or \solfalse %%%%%%%%%%%%%%%%%%%%%%%%%%%%%
\newif\ifsol
\soltrue  % %SOLUTIONS
% \solfalse % NO SOLUTIONS
%%%%%%%%%%%%%%%%%%%%%%%%%%%%%%%%%%%%%%%%%%%%%%%%%%%%%%%%%%%%%%

%%%%%%%%%%%%%%%% solution help %%%%%%%%%%%%%%%%%%%%%
\newcommand{\solu}[2]{ \begin{mdframed} \ifsol #2 \else \vspace{#1} \fi \end{mdframed} }
\newcommand{\solt}{ \ifsol (T) \fi }
\newcommand{\solf}{ \ifsol (F) \fi }
\newcommand{\solm}[1]{\ifsol \textit{(#1)} \fi}

\begin{document}
  \thispagestyle{firstpagestyle}
  \begin{center}
    {\large \textbf{Recitation 5}}
    
    {\large Induction}
  \end{center}
 
 
\section*{Induction}
 
 \subsection*{Why does induction work?}
 
 Let's consider an infinite ladder (the best kind of ladder). Suppose we can prove to you both of the following things:
 
 \begin{enumerate}
 \item You can get to the $1^{\text{st}}$ step of the ladder by stepping up to it.
 \item If you can get to the $k^{\text{th}}$ step of the ladder, then you can get to step $k+1$ by stepping up to it.
 \end{enumerate}
 
Why is it the case that for all $n \geq 1$, you can get to the $n^{\text{th}}$ step of the ladder? Discuss with your group.

\solu{2cm}{We already know we can get to the first step from the first statement. Then, we know we can get to the second step from the second statement. From there, the process repeats and we conclude that we can get to the third, then the fourth... and so on.}
 
 Why are we talking about climbing infinite ladders? Well, it turns out this is a good way to think about how induction works. 

	The \textit{base case} says that we can reach the first step of the ladder.

	The \textit{inductive hypothesis} says that we can get to the $k^{\text{th}}$ step of the ladder.

	The \textit{inductive step} says that if we can get to the $k^{\text{th}}$ step of the ladder, then we can get to step $k+1$. 

	Therefore, once we get to step 1, we can get to step 2. Once we get to step 2, we can get to step 3. And so on for all steps of the infinite ladder.
 
 \subsection*{Induction Template}
 
      We will now review the template for an inductive proof.


      For example, say we are trying to prove that $\sum_{i=0}^n i^2 = \frac{n(n+1)(2n+1)}{6}$ is true for all $n \in \N$. In other words, show that this is the equation for calculating the sum of squares $0^2 + 1^2 + 2^2 + \cdots + n^2$.

      \begin{enumerate}
        \item Define the predicate $P(n)$. Recall that a predicate is a function that takes in an argument, $n$, and evaluates to true or false.
        
         \textit{Let $P(n)$ be the predicate that $\sum_{i=0}^n i^2 = \frac{n(n+1)(2n+1)}{6}$.}
        
        \item Make the aspirational assertion that, for all $n \geq a$, where $a$ is the smallest value we are considering, $P(n)$ holds. \textbf{Remember to bound $n$}! 
        
        \textit{We will show that, for all $n \geq 0$, $P(n)$ holds.}

        \item Show that the base case is true. For some proofs, we may want multiple base cases, but not this time.

        \textit{We will first show $P(0)$ is true. $\sum_{i=0}^0 i^2 = 0$ and $\frac{0(0+1)(2*0+1)}{6} = 0$ so they are equal.}

        \item State the inductive hypothesis. In standard induction\footnote{We will also cover strong induction, in which we assume $P(i)$ is true for \textbf{all} $a \le i \leq k$}, we assume 
        
        $P(k)$ is true for some fixed, arbitrary integer $k \geq a$, where $a$ is your base case value. Sometimes, you may need multiple base cases, and you'll want $k$ to be greater than or equal to the biggest of them.
% 			If you do need multiple base cases, you'll usually detect this when writing your inductive step (so you'll need to go back and add more base cases!). In this example though, it turns out we only need one. 

        \textit{Assume $P(k)$ is true for some fixed, arbitrary integer $k \geq 0$.}

        \item Show that $P(k+1)$ is true given the inductive hypothesis. At some point, you'll want to ``invoke the inductive hypothesis", which is using the fact that $P(k)$ is true to show something else in your proof.

		\textit{We will now show that $P(k+1)$ holds, namely  $\sum_{i=0}^{k+1} i^2 = \frac{(k+1)((k+1)+1)(2(k+1)+1)}{6}$}.

       	\textit{We know that $\sum_{i=0}^{k+1} i^2 = \left( \sum_{i=0}^{k} i^2 \right)+ (k+1)^2$.}

		\textit{Invoking the inductive hypothesis, we know that $\sum_{i=0}^k i^2 = \frac{k(k+1)(2k+1)}{6}$.}

         \textit{Therefore}
		\begin{align*}
			\sum_{i=0}^{k+1} i^2 &= \left( \sum_{i=0}^{k} i^2 \right)+ (k+1)^2 \\
			&=  \frac{k(k+1)(2k+1)}{6} + (k+1)^2 \\
			&= \frac{k(k+1)(2k+1) + 6(k+1)^2}{6} \\
			&= \frac{(k+1)(k(2k+1)+6(k+1))}{6} \\
			&= \frac{(k+1)(2k^2+k+6k+6)}{6} \\
			&= \frac{(k+1)(2k^2+7k+6)}{6} \\
			&= \frac{(k+1)(k+2)(2k+3)}{6} \\
			&= \frac{(k+1)((k+1)+1)(2(k+1)+1)}{6}
		\end{align*}
		\textit{as needed.}
		
		 \item Conclude.
		 
		 \textit{Because the base case $P(0)$ holds, and because $P(k)$ IMPLIES $P(k + 1)$, we have shown by the principle of induction that for all $n \geq 0$, $P(n)$ holds.} \qed
		\end{enumerate}

\subsection*{$\star$ Note $\star$}

For the sake of time, we're only going to look for proof \textit{sketches} in recitation. It's OK to not write down everything, as long as you understand it. In your homework, we'll be looking for full-fledged formal proofs.

	\subsection*{Task 1}

      Prove by induction that, for all $n \geq 2$, 
      
      $$(1 - \frac{1}{2^2})(1 - \frac{1}{3^2}) \cdots (1 - \frac{1}{n^2}) = \frac{n+1}{2n}$$

	\solu{12cm}{
	Let P($n$) be the predicate that $$(1 - \frac{1}{2^2})(1 - \frac{1}{3^2}) \cdots (1 - \frac{1}{n^2}) = \frac{n+1}{2n}$$ defined for integers $n \geq 2$. \\
	
	\textbf{Base Case:} \\
	$LHS = 1 - \frac{1}{2^2} = \frac{3}{4}$.\\
	$RHS = \frac{2 + 1}{4} = \frac{3}{4}$. 
	\\Hence, P(2) holds. \\
	
	\textbf{Inductive Hypothesis:} Suppose that for some fixed, arbitrary integer $k \geq 2$, $P(k)$ holds. That is, $(1 - \frac{1}{2^2})\cdots(1 - \frac{1}{k^2}) = \frac{k + 1}{2k}$. \\
	
	\textbf{Inductive Step:} Now, consider $(1 - \frac{1}{2^2})\cdots(1 - \frac{1}{k^2})(1 - \frac{1}{(k + 1)^2})$.
	
	By the I.H., we know that it is equal to $\frac{k + 1}{2k}(1 - \frac{1}{(k + 1)^2})$.
	
	Distributing, this is equal to $\frac{k + 1}{2k} - \frac{k + 1}{2k(k + 1)^2}$, which simplifies into $\frac{(k + 1)^2 - 1}{2k(k + 1)}$.
	
	Expanding the numerator gives us $\frac{k^2 + 2k}{2k(k + 1)}$, and cancelling out $k$ results in $\frac{k + 2}{2(k + 1)}$, as needed. Thus, P($k+1$) is true.\\
	
	\textbf{Conclusion:} Since P(2) is true and for any integer $k \geq 2$, and $P(k)$ implies $P(k+1)$, $P(n)$ is true for all integers $n \geq$ 2.}
	
	
\subsection*{Task 2}
Blue and Green are playing a \textit{very fun game}. They have some number of tasks to do, and take turns completing one, two, or three tasks. Whoever has to do the last task loses. If Green goes second, prove that they always have a winning strategy if the number of tasks equals $4k+1$ for some $k \geq 0$.

\solu{10 cm}{\textbf{Proof}

Let us define $P(n)$ as the predicate that Green will win if the game starts with $4n+1$ tasks, for $n \geq 0$.
    
\textbf{Base case.} For $n=0$, the number of tasks is $1$ which forces Blue to perform the last task and lose.
    
\textbf{Inductive Hypothesis:} Assume $P(k)$, that Green can win if the number of tasks is $4k+1$ when $k \geq 0$.
    
\textbf{Inductive Step:} We want to prove $P(k)$ IMPLIES $P(k+1)$. The number of tasks is now $4(k+1) + 1 = 4k + 5$. Blue can choose to complete either 1, 2, or 3 tasks, making the pile $4k+4$, $4k+3$, and $4k+2$, respectively. Green can then finish 3, 2, or 1 tasks, respectively, which will make the list of size $4k + 1$. Given our inductive hypothesis, Green will win, and now has a winning strategy.
    
\textbf{Conclusion:} By induction, having proven $P(0)$ and $P(k)$ IMPLIES $P(k+1)$, if the number of tasks equals $4k+1$ for some $k$, then the second player has a winning strategy.}

	\subsection*{\textit{Optional:} Task 3}

      Use induction to prove the following
      generalization of one of De Morgan's laws:
      \begin{center}
      $\lnot(p_1 \land p_2 \land p_3 \land ... \land p_n) 
      = \lnot p_1 \lor \lnot p_2 \lor ... \lor \lnot p_n$
      \end{center}
      for $n \in \Z^+, n \geq 2$.

	\solu{17cm}{
      Claim: $\lnot(p_1 \land p_2 \land p_3 \land ... \land p_n) 
      = \lnot p_1 \lor \lnot p_2 \lor ... \lor \lnot p_n$ 
      for $n \in \Z^+, n \geq 2$.

      Define $P(n)$ as the property that 
      $\lnot(p_1 \land p_2 \land p_3 \land ... \land p_n) 
      = \lnot p_1 \lor \lnot p_2 \lor ... \lor \lnot p_n$
      for $n$ predicates.

      \textbf{Base case:} We prove $P(2)$.

      $\lnot(p_1 \land p_2) = \lnot p_1 \lor \lnot p_2$. This is the simple form
      of one of De Morgan's Laws that has been previously proven. So, $P(2)$ is
      true.

      \textbf{Inductive hypothesis:} 
      Assume $P(k)$ is true for some $k \in \Z^+$,
      $k \geq 2$. That is, 
      $\lnot(p_1 \land p_2 \land p_3 \land ... \land p_k) 
      = \lnot p_1 \lor \lnot p_2 \lor ... \lor \lnot p_k$.

      \textbf{Inductive step:} We will now prove $P(k + 1)$. 
      That is, we will show that
      $\lnot(p_1 \land p_2 \land p_3 \land ... \land p_k \land p_{k+1}) 
      = \lnot p_1 \lor \lnot p_2 \lor ... \lor \lnot p_k \lor \lnot p_{k+1}$.

      Let $q_k = p_1 \land p_2 \land \dots \land p_k$.

      Then, $\lnot(p_1 \land p_2 \land p_3 \land ... \land p_k 
      \land p_{k+1}) = \lnot(q_k \land p_{k+1})$.

      By Simple De Morgan's law, 
      $\lnot(q_k \land p_{k+1}) = \lnot q_k \lor \lnot p_{k+1}$.

      By the inductive hypothesis, $\lnot q_k
      = \lnot(p_1 \land p_2 \land \dots \land p_k)
      = \lnot p_1 \lor \lnot p_2 \lor \dots \lor \lnot p_k$.

      By substituting $q_k$ back in, we have that
      $\lnot(p_1 \land p_2 \land p_3 \land ... \land p_k \land p_{k+1}) 
      = \lnot p_1 \lor \lnot p_2 \lor ... \lor \lnot p_k \lor \lnot p_{k+1}$.

      Therefore, $P(k)$ implies $P(k+1)$.

      Because the base case and inductive step are true, $P(n)$ is true for all
      $n\in Z^+$, $n\geq 2$ by induction.

    }
\textbf{Checkpoint 1 — Call over a TA!}

\newpage

 \section*{Strong Induction}
 
With standard induction under our belts, it's time to look at a variant of it, strong induction. In many ways, strong induction is similar to normal induction, as the basic steps listed above are all the same.

Remember that our goal is to prove $\forall i\geq a,  P(i)$.
The difference is in the inductive hypothesis. When using induction, we assume that $P(k)$ is true to prove $P(k+1)$. In strong induction, we assume that the particular statement holds at all the steps from the base case to the $k$th step. Sometimes, we assume all of $P(b)$, $P(b+1)$, ..., $P(k)$ are true to prove $P(k+1)$, where $b$ is the base case. Note that we may need to prove the base case for multiple values. 

Why would we need to do that? Sometimes, you can't just rely on the fact that $P(k)$ is true. Maybe you also need $P(k-1)$ to be true, or perhaps also $P(k-2)$, or even $P(k/2)$. While writing out your inductive step, if you realize that $P(k)$ isn't enough to prove $P(k+1)$, odds are you need strong induction.

\section*{Multiple Base Cases?}

Something to think about: If you need both $P(k)$ and $P(k-1)$, you also need \textit{multiple base cases}. Say you're trying to prove $P(n)$ for $n \geq 1$. If you prove $P(1)$ as your base case, how can you show $P(2)$ without $P(0)$ in the inductive step? You'd have to include \textit{both} $P(1)$ and $P(2)$ as base cases.

And so, you naturally might ask...

\subsection*{Which method should we use?}

With some standard types of problems (e.g., sum formulas) it is clear ahead of time what type of induction is likely to be required, but usually this question answers itself during the exploratory/scratch phase of the argument. In the induction step you will need to reach the $k + 1$ case, and you should ask yourself which of the previous cases you need to get there. If all you need to prove the $k + 1$ case is the case $k$ of the statement, then ordinary induction is appropriate. On the other hand, you may realize that you need the two preceding cases ($k - 1$ and $k$) or the full range of preceding cases, to get to $k + 1$, in which case strong induction is needed.

\section*{Example}

Prove that every integer $n \geq 2$ can be written as a product of one or more prime numbers.

\textbf{Proof}

Let $P(n)$ be the predicate “$n$ can be written as a product of one or more prime numbers”.

\textbf{Base case.} The integer 2 is prime, so it is a product of exactly one prime number (itself). Therefore, $P(2)$ is true.

\textbf{Inductive Hypothesis.} Assume the inductive hypothesis, that for a particular $k$, $P(i)$ is true for all $2 \leq i \leq k$.




\textbf{Inductive Step.} We must prove $P(k + 1)$, that $k + 1$ is the product of one or more prime numbers. $k + 1$ is either prime or composite. If it is prime, then it is the product of exactly one prime number (itself), and $P(k + 1)$ is true. If it is composite, then by definition it is the product of two factors, $k + 1 = ab$, where $a$ and $b$ are integers $\geq 2$.

Since $a$ and $b$ are both greater than 1, they must also both be less than $k + 1$. By the inductive hypothesis, $a$ and $b$ can each by written as a product of one or more primes. But since $k + 1 = ab$, we can combine these two products to express $k + 1$ as a product of primes, so $P(k + 1)$ is true.


\smallskip 

Thus \textbf{inductive hypothesis and inductive step} imply:

$$\forall k, (\bigwedge_{j=2}^k P(j))\implies P(k+1)$$

\textbf{Conclusion.} Since $P(2)$ is true and $P(2),...,P(k)$ together imply $P(k + 1)$, $P(n)$ is true for all integers $n \geq 2$.

\newpage
\subsection*{Task 4}

\begin{enumerate}
    \item In the above example proof, why did we only need one base case?
    
    \solu{1cm}{We had the guarantee that $a$ and $b$ are between 2 and $k$, so there was no risk of them being less than our highest base case.}
    
    \item In which step of the induction proof is there a difference between ordinary and strong induction? What is the difference? 
    
    \solu{1cm}{Inductive Hypothesis! (see description of strong induction above). In weak induction, we only assume that particular statement holds at k-th step, while in strong induction, we assume that the particular statement holds at all the steps from the base case to the k-th step.

    Note: A common incorrect answer here is 'multiple base cases for strong induction
'.}
\end{enumerate}

\subsection*{Task 5}

\begin{enumerate}
    \item Prove by strong induction that every amount of postage that is at least 15 cents can be made from 4-cent and 5-cent stamps.

\textit{Hint:} You'll need 4 base cases.

\solu{10 cm}{\textbf{Proof}

Let $P(n)$ be the predicate that a $n$-cent postage, where $n \geq 15$, can be made from 4-cent and 5-cent stamps.

\textbf{Base case.} If the postage is 15 cents, we can make it with three 5-cent stamps. If the postage is 16 cents, we can make it with four 4-cent stamps. If it is 17, we can use three 4-cent stamps plus one 5-cent stamps. If it is 18, we use two 4-cent stamps and two 5-cent stamps.

\textbf{Inductive Hypothesis.} Assume the inductive hypothesis, that for a particular $k$, $P(i)$ is true for all $15 \leq i \leq k$, where $k$ is an integer with $k \geq 18$; in other words, suppose that we have shown how to construct postages for every value from 15 up through $k$.

\textbf{Inductive Step.} We must prove $P(k + 1)$, that a postage of $k + 1$ cents can be made from 4-cent and 5-cent stamps. Since we have already proved the base cases up through 18 cents, we will assume that $k + 1 \geq 19$. Since $k + 1 \geq 19$, $(k + 1) - 4 \geq 15$. So, by the inductive hypothesis, we can construct postage for $(k + 1) - 4$ cents using $m$ 4-cent stamps and $n$ 5-cent stamps, for some natural numbers $m$ and $n$. In other words, $(k + 1) - 4) = 4m + 5n$. But then $k + 1 = 4(m + 1) + 5n$. So we can construct $k + 1$ cents of postage using $m + 1$ 4-cent stamps and $n$ 5-cent stamps, which is what we needed to show.

\textbf{Conclusion.} Since $P(15)$, $P(16)$, $P(17)$, and $P(18)$ is true and $P(15),...,P(k)$ together imply $P(k + 1)$, $P(n)$ is true for all integers $n \geq 15$.

Notice that we needed to directly prove four base cases, since we needed to reach back four integers in our inductive step. It’s not always obvious how many base cases are needed until you work out the details of your inductive step.\footnote{https://courses.engr.illinois.edu/cs173/sp2009/Lectures/lect\_18.pdf}.}

\item Why do we need to use strong induction for this proof? In other words, why can't we just 'assume that $P(k)$ is true' in our induction hypothesis? 

\solu{1cm}{Our proof relies on the observation that given postage for $k+1$ cents, we can set aside a 4-cent stamp and be left with $k+1-4=k-3$ cents of postage to make. As such, we assume that $P(k-3)$ is true in our inductive step, so using an inductive hypothesis $P(k) \implies P(k+1)$ is not sufficient, and instead we use strong induction!}

\item In part 1, we proved by strong induction that every amount of postage that is at least 15 cents can be made from 4-cent and 5-cent stamps.

Prove the same claim using \textbf{ordinary induction}.

\solu{10 cm}{\textbf{Proof}

Let $P(n)$ be the predicate that a $n$-cent postage, where $n \geq 15$, can be made from 4-cent and 5-cent stamps.

\textbf{Base case.} If the postage is 15 cents, we can make it with three 5-cent stamps.

\textbf{Inductive Hypothesis.} Assume the inductive hypothesis, that for a particular integer $k \geq 12$, $P(k)$ is true. In other words, we assume that a postage of $k$ cents can be made from 4-cent and 5-cent stamps.

\textbf{Inductive Step.} We must prove $P(k + 1)$, that a postage of $k + 1$ cents can be made from 4-cent and 5-cent stamps.

By the inductive hypothesis, there is a way to make postage of $k$ cents using only 4-cent and 5-cent stamps. We can then make postage for $k+1$ cents in one of two ways:

\begin{itemize}
    \item Case 1: The postage for $k$ cents includes a 4-cent stamp. In this case, we can make $k+1$ cents of postage by removing the 4-cent stamp and replacing it with a 5-cent stamp, which will increase the overall amount by 1 cent. In other words, we get $k-4+5=k+1$ cents of package.
    \item Case 2: The postage for $k$ cents does not include a 4-cent stamp. In this case, the postage for $k$ cents must contain at least three 5-cent stamps (since $k \geq 15$). We can then make $k+1$ cents of postage by removing the three 5-cent stamps and replacing them with four 4-cent stamps, which will increase the overall amount by 1-cent. In other words, we get $k-3 \cdot 5 + 4 \cdot 4 = k+1$ cents of package.
\end{itemize}

Since one of these cases must occur, then we can make $k+1$ cents of postage, which shows that $P(k+1)$ is true.


\textbf{Conclusion.} Because $P(1)$ is true, and because, for all $k$, $P(k)$ implies $P(k+1)$, we conclude that $P(n)$ is true for all integers $n \geq 1$.}

\item \textit{Optional: } Are strong and ordinary induction equivalent?

\solu{2cm}{Yes, strong and ordinary induction are logically equivalent! See \href{http://mathcenter.oxford.emory.edu/site/math125/strongInductionEquivalence/}{\color{blue} this website} for the proof of their equivalence.

However, sometimes it may be easier to use one method over the other. The section "Which method should we use?" above also discusses this briefly.}
\end{enumerate}

\subsection*{\textit{Optional:} Task 6}

Consider a candy bar with $n$ squares in a row. Suppose we want to break this candy bar up into individual squares. How many breaks should we perform? 
	
\begin{figure}[tph!]
\centerline{\includegraphics[totalheight=3cm]{chocolate.jpg}}
    \caption{A Delicious Candy Bar}
    \label{fig:verticalcell}
\end{figure}
	
\textbf{Claim}: For all $n \geq 1$, \textbf{any sequence} of $n-1$ breaks will reduce a candy bar of $n$ squares into single squares. This means it doesn't matter what order we break squares in: think about this fact for your inductive step!

Prove this claim by (strong) induction. 

\solu{9cm}{Let $P(n)$ be the predicate that, given a bar of $n$ squares, any sequence of $n-1$ breaks will reduce it to single squares. 

\textbf{Base Case:} $P(1)$: Given a bar with 1 square, 0 breaks leaves us with one square. Hurrah!

\textbf{Inductive Hypothesis:} Let $k$ be some arbitrary positive integer. For all positive integers $i \leq k$, assume $P(i)$. That is, assume that for a bar of length $i$, any sequence of $i-1$ breaks will reduce it to single squares.

\textbf{Inductive Step:}  Consider a candy bar with $k + 1$ squares.
Make some break. Let's say we make it after the $p$th square, where $p$ is some positive integer. Then, we have two candy bars, one of length $p$, the other of length $k + 1 - p$. Both $p$ and $k + 1 - p$ are positive integers less than $k+1$, as we are breaking the bar into two positive pieces. By our inductive hypothesis, we know that it takes $p - 1$ breaks to break up the first, and $(k + 1) - p - 1$ to break up the second. These breaks result in a total of: $(p - 1) + ((k + 1 - p) - 1) = k + 1 - 2$ breaks. We then have that our initial break plus this number of breaks is: $1 + (k + 1 - 2) = (k + 1) - 1$ breaks, as needed. That shows $P(k+1)$ is true. \\

\textbf{Conclusion}: As $P(1)$ is true and for any positive integer $k$, and $P(1), \ldots. P(k)$ imply $P(k+1)$, $P(n)$ is true for all positive integers $n$.}


\subsection*{Required: Feedback Form}
Congrats on finishing Recitation~5---it's remarkable how much we learned about proof techniques, logic, set theory, functions and relations, and induction!

The last part of recitation is to fill out this mandatory \href{https://forms.gle/i9Zkq99XyaoweysF8}{\color{blue} feedback form}, which should be completed individually. You will be checked off only if we receive a submission from you.

Your responses will be super helpful as we continue to improve the course, and thank you for an amazing first few weeks of 22!

\subsection*{Checkoff}
If you are done, call over a TA to get checked off. There's a bonus problem on the next page while you wait.

Lastly, just a reminder that you can direct conceptual questions to your TAs during recitation as well!
\newpage

\subsection*{\textit{Optional Challenge:} Green-Eyed Aliens}

There are $n$ (where $n$ is 2 or greater) aliens sitting in a circle so that every alien can see every other alien.

Every alien has green eyes. However, no alien knows its own eye color. Additionally, the aliens cannot talk, so they cannot inform each other of the fact that they have green eyes. However, if an alien ever figures out their own eyes are green, they will leave the spaceship that night.

On Day 1, Professor Littman comes and tells the circle of aliens that at least one of them has green eyes. Prove that, on the $n$th night, all $n$ aliens will leave the ship. 

\solu{6cm}{Let $P(n)$ be the predicate that $n$ aliens will leave on the $n$th night.

\textbf{Base Case:} If there are 2 aliens, then it takes them two nights to figure out they both have green eyes. To arrive at this conclusion, consider one of the aliens, $D$. On the first day, $D$ sees that the other alien, $D'$ has green eyes. $D'$ does not know that $D'$ has green eyes, but they see that $D$ does. Hence, neither $D$ nor $D'$ realizes on the first day that they themselves have green eyes. On the second day though, $D$ sees that the other alien, $D'$, did NOT realize that $D'$ has green eyes. $D'$ would have come to this conclusion if $D'$ could see $D$'s eyes were not green given Professor Littman's claim that at least one of them has green eyes. Hence, $D$ knows that $D$ has green eyes! $D'$ also knows that $D'$ has green eyes by the same reasoning! (These aliens are strong at logic.) So, on the second night, they both leave. Thus, $P(2)$ is true.

\textbf{Inductive Hypothesis:} Now, for some arbitrary integer $k \geq 2$, assume $P(k)$. That is, it takes $k$ nights for $k$ aliens to leave.

\textbf{Inductive Step:} Consider the case where there are $k+1$ aliens and consider what one individual alien is thinking. This one alien sees $k$ aliens with green eyes. Therefore, if this alien does not have green eyes, they would expect the other aliens to leave after $k$ nights since they can see them, by the IH. Since they don't, this alien realizes that their eyes are also green and leaves on night $k+1$. Thus, $P(k+1)$ is true.

\textbf{Conclusion:} As $P(2)$ is true and, for any integer $k \geq 2$ $P(k)$ implies $P(k+1)$, $P(n)$ holds for all integers $n \geq 2$.}

\end{document}